\documentclass[11pt,a4paper]{article}
\usepackage[utf8]{inputenc}
\usepackage[T1]{fontenc}
\usepackage[turkish]{babel}
\usepackage[margin=2.3cm]{geometry}
\usepackage{graphicx}
\usepackage{booktabs}
\usepackage{array}
\usepackage{hyperref}
\usepackage{longtable}
\title{\textbf{BLM308 Veri Madenciliği Final Proje Raporu}\\Banka Pazarlama Kampanyalarında Tahmin Edilebilirlik ve Açıklanabilirlik}
\author{Sude ANLAŞ -- 231041038 \and Doğan Fatih OĞUR -- 231041048 \and Merve Naz ORAN -- 231041058}
\date{Bahar 2026}
\begin{document}
\maketitle
\tableofcontents
\newpage

\section*{Takım Bilgileri ve Rol Dağılımı}
\begin{longtable}{p{3cm}p{2.3cm}p{4cm}p{6cm}}
\toprule
Ad Soyad & Numara & Rol & Yaptığı/Yazdığı Bölümler\\
\midrule
Sude ANLAŞ & 231041038 & Veri Mühendisi & Veri seti, EDA, ön işleme, veri sızıntısı analizi\\
Doğan Fatih OĞUR & 231041048 & Model Mühendisi & Random Forest/MLP, Streamlit arayüzü, kod düzeni\\
Merve Naz ORAN & 231041058 & Değerlendirme ve XAI Analisti & Metrikler, McNemar testi, SHAP-LIME yorumları, sonuç ve iş değeri analizi\\
\bottomrule
\end{longtable}

\section*{Özet}
Bu çalışmada UCI Bank Marketing veri seti kullanılarak müşterinin vadeli mevduata abone olup olmayacağı tahmin edilmiştir. Proje CRISP-DM adımlarına göre tasarlanmış; veri anlayışı, ön işleme, modelleme, değerlendirme ve iş değeri yorumunu birlikte ele almıştır. Logistic Regression ve Decision Tree temel modeller, Random Forest ve MLP ise black-box modeller olarak kullanılmıştır. Model başarısı accuracy, precision, recall, F1 ve ROC-AUC metrikleriyle ölçülmüştür. Ayrıca kampanya sonrası bilinen \texttt{duration} değişkeninin veri sızıntısı riski tartışılmış, bu nedenle \texttt{duration} dahil ve hariç iki senaryo karşılaştırılmıştır. Random Forest kararları SHAP ve LIME ile açıklanarak genel ve yerel model davranışı yorumlanmıştır.

\textbf{Anahtar kelimeler:} veri madenciliği, sınıflandırma, Random Forest, MLP, SHAP, LIME, XAI.

\section{Giriş}
Bankalar pazarlama kampanyalarında sınırlı çağrı merkezi kapasitesini doğru müşterilere yönlendirmek ister. Rastgele arama stratejileri hem maliyeti artırır hem de müşteri memnuniyetini düşürebilir. Bu projenin problemi, geçmiş kampanya ve müşteri bilgilerinden yararlanarak bir müşterinin vadeli mevduata abone olma olasılığını tahmin etmek ve black-box model kararlarını açıklanabilir hale getirmektir.

Bu problemin çözümü bankanın daha verimli müşteri hedeflemesi yapmasını sağlar. Ancak finans alanında kararların gerekçelendirilmesi gerekir; bu nedenle XAI bölümü projenin merkezindedir. Pazarlama yöneticisi yalnızca model skorunu değil, bu skora hangi değişkenlerin neden olduğunu da görebilmelidir.

\section{Yöntem}
Proje CRISP-DM'in altı adımına göre yürütülmüştür: iş anlayışı, veri anlayışı, veri hazırlığı, modelleme, değerlendirme ve dağıtım/iş değeri yorumu. Python, scikit-learn, SHAP ve LIME kullanılmıştır. Weka karşılıkları olarak J48, RandomForest ve MultilayerPerceptron algoritmaları raporlanabilir durumdadır.

\section{Veri}
Veri seti UCI Machine Learning Repository üzerinde yayımlanan Bank Marketing veri setidir. \texttt{bank-additional-full.csv} dosyası 41.188 gözlem ve 20 girdi değişkeni içerir. Hedef değişken \texttt{y}, müşterinin vadeli mevduata abone olup olmadığını gösteren ikili sınıftır.

\begin{table}[h]
\centering
\begin{tabular}{ll}
\toprule
Özellik & Değer\\
\midrule
Örnek sayısı & 41.188\\
Girdi değişkeni & 20\\
Hedef değişken & y: yes/no\\
Pozitif sınıf & yes = 4.640\\
Negatif sınıf & no = 36.548\\
\bottomrule
\end{tabular}
\caption{Veri seti özeti}
\end{table}

Kategorik değişkenler One-Hot Encoding ile dönüştürülmüş, sayısal değişkenler StandardScaler ile ölçeklendirilmiştir. Hedef değişken \texttt{yes=1}, \texttt{no=0} olarak kodlanmıştır. Veri stratified yöntemle yüzde 75 eğitim, yüzde 25 test olacak şekilde ayrılmıştır.

\section{Modelleme}
Seçilen modeller hem temel yorumlanabilir yaklaşımları hem de XAI gerektiren black-box yaklaşımları kapsar.

\begin{table}[h]
\centering
\begin{tabular}{llll}
\toprule
Model & Tür & Weka Karşılığı & Rol\\
\midrule
Logistic Regression & Lineer & Logistic & Temel model\\
Decision Tree & Ağaç & J48 & Açıklanabilir karşılaştırma\\
Random Forest & Ensemble & trees.RandomForest & Ana black-box model\\
MLP & Sinir ağı & functions.MultilayerPerceptron & Alternatif black-box model\\
\bottomrule
\end{tabular}
\caption{Model ailesi}
\end{table}

\section{Sonuçlar}
\begin{table}[h]
\centering
\small
\begin{tabular}{llrrrrr}
\toprule
Senaryo & Model & Accuracy & Precision & Recall & F1 & ROC-AUC\\
\midrule
duration dahil & Logistic Regression & 0.865 & 0.450 & 0.909 & 0.602 & 0.943\\
duration dahil & Decision Tree & 0.845 & 0.415 & 0.910 & 0.570 & 0.921\\
duration dahil & Random Forest & 0.868 & 0.457 & 0.925 & 0.611 & 0.949\\
duration dahil & MLP & 0.915 & 0.650 & 0.528 & 0.582 & 0.947\\
duration hariç & Logistic Regression & 0.835 & 0.367 & 0.644 & 0.467 & 0.805\\
duration hariç & Decision Tree & 0.822 & 0.344 & 0.635 & 0.446 & 0.773\\
duration hariç & Random Forest & 0.864 & 0.430 & 0.633 & 0.512 & 0.816\\
duration hariç & MLP & 0.900 & 0.647 & 0.245 & 0.355 & 0.808\\
\bottomrule
\end{tabular}
\caption{Model performans sonuçları}
\end{table}

\texttt{duration} dahil senaryoda Random Forest en yüksek ROC-AUC değerine ulaşmıştır. Ancak \texttt{duration} çağrı tamamlandıktan sonra bilindiği için kampanya öncesi müşteri seçimi açısından veri sızıntısı riski taşır. Bu nedenle gerçek kullanımda \texttt{duration} hariç senaryo daha dürüst değerlendirme olarak kabul edilmelidir.

\section{XAI: SHAP ve LIME}
XAI bölümünün amacı black-box modelin yalnızca skor üretmesini değil, karar gerekçesini de okunabilir hale getirmektir. SHAP global düzeyde hangi değişkenlerin model kararlarını daha fazla etkilediğini gösterir. LIME ise tek bir müşteri çevresinde yerel ve yorumlanabilir bir açıklama modeli kurar.

Random Forest önem sıralaması, ekonomik göstergeler, önceki kampanya sonucu ve kampanya etkileşim değişkenlerinin model kararında belirgin rol oynadığını göstermiştir. LIME çıktısı ise tek müşteri düzeyinde hangi koşulların \texttt{yes} sınıfını desteklediğini veya zayıflattığını göstermiştir.

\section{İş Değeri ve Sınırlılıklar}
Model, kampanya listesinde müşterileri olasılık skoruna göre sıralamak için kullanılabilir. En yüksek olasılıklı müşterilere öncelik verilmesi çağrı merkezi kaynaklarının daha verimli kullanılmasını sağlar. XAI çıktıları ise pazarlama ekibine segment bazlı içgörü sunar.

Sınırlılıklar arasında veri setinin tek ülke ve belirli dönem kampanyalarından oluşması, sınıf dengesizliği ve \texttt{duration} değişkeni gibi süreç sonrası oluşan değişkenlerin kullanımı yer alır. Daha ileri çalışmalarda maliyet duyarlı öğrenme, eşik optimizasyonu ve zaman bazlı validasyon denenebilir.

\section{Sonuç}
Bu proje, Bank Marketing veri seti üzerinde black-box sınıflandırma modelleri kurmuş ve bu modellerin kararlarını SHAP/LIME ile açıklamıştır. Random Forest performans ve açıklanabilirlik dengesinde güçlü bir aday olarak öne çıkmıştır. Proje, veri madenciliği sürecinde teknik metrikler kadar veri hazırlığı, veri sızıntısı farkındalığı, XAI ve iş değeri yorumunun da önemli olduğunu göstermektedir.

\section*{Kaynaklar}
\begin{enumerate}
\item Moro, S., Cortez, P., Rita, P. (2014). A Data-Driven Approach to Predict the Success of Bank Telemarketing. Decision Support Systems.
\item UCI Machine Learning Repository. Bank Marketing Dataset. \url{https://archive.ics.uci.edu/dataset/222/bank+marketing}
\item Lundberg, S. M., Lee, S.-I. (2017). A Unified Approach to Interpreting Model Predictions. NeurIPS.
\item Ribeiro, M. T., Singh, S., Guestrin, C. (2016). Why Should I Trust You? Explaining the Predictions of Any Classifier. NAACL-HLT.
\item Witten, I. H., Frank, E., Hall, M. A., Pal, C. J. Data Mining: Practical Machine Learning Tools and Techniques.
\end{enumerate}

\section*{Yapay Zeka Kullanımı Beyanı}
Bu raporun taslak oluşturma ve dil düzenleme aşamasında yapay zeka desteği kullanılmıştır. Veri seti, modelleme kodu, grafikler ve sonuç yorumları proje kapsamında kontrol edilerek düzenlenmiştir.

\end{document}
